\begin{frame}{비교 프로토콜 (1/3): 시드 $\geq 3$ --- 왜인가}
  \begin{itemize}
    \item ``한 번의 실행''은 아무것도 말해주지 않는다:
    \begin{itemize}
      \item $\varepsilon$-greedy 탐험의 무작위성,
      \item (FrozenLake처럼) \textbf{환경 전이 자체}의 무작위성,
      \item 정책 초기화.
    \end{itemize}
    \item $\to$ \textbf{같은 코드·같은 파라미터를 그대로 돌려도 매번 다른
      학습 곡선}이 나온다.
    \item \kb{시드를 바꾸는 것}은 ``실험을 다시 하는 것''이 아니라
      \textbf{``우연에 좌우되는지 재검사하는 것''}.
  \end{itemize}
  \vskip0.6em
  \begin{block}{프로토콜의 세 규칙}
    \begin{enumerate}
      \item \textbf{시드 $\geq 3$} (예: 0, 1, 2)으로 알고리즘을 \textbf{전부
        다시} 학습.
      \item \textbf{학습 도중 성과}와 \textbf{학습 후 평가 성과}를 반드시
        \textbf{분리}해 보고.
      \item 비교는 \textbf{시드별 숫자 자체}(예: ``$-13$ / $-13$ / $-13$'')를
        표로 제시하고, 평균$\pm$표준편차나 ``일관된 경향''으로 결론.
    \end{enumerate}
  \end{block}
\end{frame}
